\section{Introduction}
\label{sec:intro}

A turbine blade cracks after 12,000 flight hours. A bearing seizes in a wind turbine gearbox during peak generation. A pump seal fails in a chemical plant, triggering an emergency shutdown. These are \emph{grey swans}---events that are rare in operational data but follow known degradation physics and are, in principle, predictable~\citep{greyswan_supplychain2025}. Unlike black swans, which are fundamentally unforeseeable~\citep{taleb2007blackswan}, grey swans are physically plausible failures that sit in the long tail of operating experience. They present a data paradox for prognostics: \emph{the events most critical to predict are precisely the ones with the least training data}. A bearing that runs for 10,000 hours before failing provides abundant healthy-state data but perhaps a single failure trajectory.

Modern predictive maintenance promises to catch these failures before they cascade into catastrophic downtime. Deep learning approaches for remaining useful life (RUL) prediction have shown strong results on benchmark datasets~\citep{lei2018machinery, zhao2019deeplearningphm}, with architectures ranging from CNN-LSTMs to Transformer-based models achieving impressive in-domain accuracy. However, these methods share two critical limitations. First, they require extensive labeled failure data---run-to-failure trajectories where degradation is observed from inception to failure---which is expensive, dangerous, and scarce by definition. Second, they rely on handcrafted spectral features (spectral centroid, kurtosis, envelope analysis) designed by domain experts, limiting applicability to well-studied component types.

Meanwhile, time series foundation models have emerged as a paradigm for general-purpose temporal representation learning. Models such as TimesFM~\citep{das2024timesfm}, Chronos~\citep{ansari2024chronos}, and MOMENT~\citep{goswami2024moment} demonstrate strong zero-shot forecasting by pretraining on massive time series corpora. Yet these models solve a different problem: they forecast the \emph{next value} in a sequence. For mechanical prognostics, predicting the next vibration sample is neither necessary nor sufficient---a perfect one-step-ahead predictor of raw waveform amplitude tells us nothing about remaining useful life. What matters is detecting the \emph{spectral signature changes} that track degradation: the shift in spectral centroid as surface roughness increases, the redistribution of vibrational energy across frequency bands as fatigue cracks propagate. This distinction---\emph{representing what matters for degradation} rather than reconstructing everything---motivates a targeted self-supervised approach designed for the structure of mechanical failure data.

We propose a self-supervised framework that addresses this data paradox by pretraining on abundant unlabeled operational data and requiring only a small number of failure trajectories for fine-tuning. The framework has three components:

\begin{enumerate}[nosep]
  \item \textbf{JEPA pretraining} learns waveform texture and periodicity from unlabeled vibration data across multiple bearing datasets, using masked patch prediction in a joint embedding predictive architecture~\citep{assran2023ijepa}. The encoder learns to represent vibration signals without any degradation labels.

  \item \textbf{Temporal contrastive learning} exploits the natural ordering of run-to-failure episodes to learn what \emph{changes} over a component's life. By training the encoder to distinguish early-life from late-life snapshots, it captures the spectral centroid shift that is the dominant degradation indicator for bearings.

  \item \textbf{Hybrid fusion} combines self-supervised embeddings with domain-informed spectral features. Rather than treating self-supervised and handcrafted approaches as competing, we show they capture complementary information---JEPA learns fine-grained waveform features while spectral statistics capture the monotonic degradation trend.
\end{enumerate}

\plannedc{We further extend this framework with multivariate input encoding that exploits physics-aware sensor groupings (vibration, current, thermal) through spatiotemporal masking, and a synthetic grey swan augmentation strategy that generates physics-informed failure trajectories for rare degradation modes.}

A key finding is that different self-supervised objectives learn fundamentally different representations. JEPA's masked prediction objective captures waveform texture (PC1 correlation with spectral centroid: $0.071$), while temporal contrastive learning captures spectral centroid shift (PC1 correlation: $0.856$). This mechanistic understanding explains a practical trade-off: JEPA excels within-dataset where fine-grained waveform features transfer, while contrastive learning dominates cross-dataset transfer where only invariant degradation dynamics generalize.

\paragraph{Contributions.}

\begin{enumerate}[nosep]
  \item A self-supervised framework for mechanical prognostics that learns degradation-aware representations from unlabeled vibration data across 8 bearing datasets (33,939 pretraining windows). JEPA pretraining alone matches expert-designed handcrafted features ($p = 0.40$); their fusion achieves RMSE $0.055 \pm 0.004$---a $20.7\%$ improvement over the best handcrafted-only method (\cref{sec:indomain}).

  \item Cross-dataset RUL transfer via temporal contrastive learning, achieving $+38\%$ over time-only baselines ($p < 0.001$) on FEMTO$\to$XJTU-SY---the first such result for bearing prognostics (\cref{sec:transfer}).

  \item A mechanistic analysis showing that JEPA captures waveform texture while contrastive learning captures spectral centroid shift, explaining when and why each objective excels (\cref{sec:encoder_analysis}).

  \plannedc{\item Extensions to multivariate inputs with physics-aware spatiotemporal masking (\cref{sec:multivariate}), validation on C-MAPSS turbofan engines (\cref{sec:cmapss}), and synthetic grey swan augmentation for rare failure modes (\cref{sec:augmentation_exp}).}
\end{enumerate}
